% !TeX program = xelatex
\documentclass[11pt,a4paper]{ctexart}
\usepackage[margin=26mm,headheight=15pt]{geometry}
\usepackage{amsmath,amssymb,amsthm,mathtools,booktabs,array}
\usepackage{enumitem,fancyhdr,xcolor}
\usepackage[unicode,colorlinks=true,linkcolor=blue!45!black,citecolor=blue!45!black,urlcolor=blue!45!black]{hyperref}
\newtheoremstyle{cnplain}{7pt}{7pt}{\normalfont}{}{\bfseries}{。}{0.5em}{}
\theoremstyle{cnplain}
\newtheorem{theorem}{定理}[section]
\newtheorem{proposition}[theorem]{命题}
\newtheorem{lemma}[theorem]{引理}
\newtheorem{corollary}[theorem]{推论}
\theoremstyle{definition}
\newtheorem{definition}[theorem]{定义}
\newtheorem{example}[theorem]{例}
\theoremstyle{remark}
\newtheorem{remark}[theorem]{注}
\renewcommand{\proofname}{证明}
\newcommand{\Q}{\mathbb Q}
\newcommand{\C}{\mathbb C}
\newcommand{\dR}{\mathrm{dR}}
\newcommand{\id}{\operatorname{id}}
\newcommand{\im}{\operatorname{im}}
\newcommand{\End}{\operatorname{End}}
\newcommand{\Hom}{\operatorname{Hom}}
\newcommand{\Vis}{\operatorname{Vis}}
\newcommand{\soc}{\operatorname{soc}}
\newcommand{\head}{\operatorname{head}}
\newcommand{\rank}{\operatorname{rank}}
\setlength{\emergencystretch}{2em}
\pagestyle{fancy}
\fancyhf{}
\fancyhead[L]{\small 费曼周期：操作、最小化与实际动机}
\fancyhead[R]{\small 理论补篇}
\fancyfoot[C]{\thepage}
\title{从操作保护到实际动机可见性\\[5pt]
\large 有限图表的相容最小化与秩八 banana 的完整双标架分类}
\author{费曼周期研究项目\quad 理论证明稿}
\date{2026年10月13日}
\begin{document}
\maketitle
\begin{abstract}
操作表示中的保护判据只有在明确的范畴比较下，才能用于实际费曼动机。
本文将这个接口拆成三个可证明的问题：已给定实际态射的有限图表如何同时最小化；
一个新的 Betti 操作何时能下降为最小商上的动机态射；
固定链而改变分子读出时，哪些纯核和扩张必须共同保留。
我们给出实际图表的有限长度闭包、商上的余模相容判据，以及 head 与 socle 的互补作用。
对于本项目已建立的指定秩八 banana 对象，进一步完整列出全部子对象，
从而对任意有理 Betti 链和有理 de Rham 读出给出最小对象表。
原物理链的顶层分量非零时，最小秩恰分为 $8,2,1,0$，相应非零周期的幺幂深度为 $2,1,0$。
纯秩六最大割来自链的改变，不能由保留原链的读出选择单独实现。
证明仅使用已有实际对象与群结构，不新增数值积分或验证计算；
也不把形式周期结论解释为未经证明的复数周期独立性。
\end{abstract}
\tableofcontents

\section{问题与输入：先确定操作所在的范畴}
\subsection{三个不同的对象}
设一个费曼问题已给出相对几何、积分链、形式和减除方案。
同一研究中经常同时出现三个结构：由函数解析延拓生成的单值表示；
由有限矩阵及允许操作构成的线性系统；承载 Betti、de Rham 和 Hodge 实现的实际动机。
它们都可以使用核、像、循环生成和矩阵系数语言，但其子对象与态射一般不同。

原稿 \nolinkurl{protected_period_network} 的 head 判据作用于已正确构造的单值模。
原稿 \nolinkurl{framed_exact_classification} 的 Hankel 最小化作用于已提交的有限操作表示。
这两个范围均应保留。本文的任务不是将原矩阵重新命名为动机，
而是给出实际对象与操作之间所需的闭包、下降及标架条件。

最小矩阵系数对象的存在与唯一性是既有理论，见 Brown\cite[\S2.4]{brown}；
中性 Tannaka 等价及其满忠实性见 Deligne--Milne\cite[Theorem 2.11]{dm}。
本文在这些基础上直接证明项目所需的接口和指定对象分类，不宣称上述基础理论的优先权。

\subsection{约定}
固定特征零域 $k$ 上的中性 Tannakian 范畴 $\mathcal C$，
以及忠实正合的 Betti 纤维函子 $\omega_B$。
另有相容的 de Rham 纤维函子 $\omega_{\dR}$；为简洁，两者暂取值于同一域 $k$。
对于本文指定的 $\Q$ 纤维，两种实现均为 $\Q$ 向量空间。
对更一般数域动机，若 de Rham 系数扩张，下文的子对象仍须在原范畴中取，
不能省略下降条件，也不能直接将扩域后的不变子空间视为原有理子对象。

以同调方向记对象 $E$、链 $\sigma\in E_B$ 和读出
$\lambda\in E_{\dR}^{\vee}$。标架周期的上同调对象则是 $E^\vee$，
二者的秩相同。对一个实际子对象 $J\subset E$，
``$\lambda$ 消去 $J$''是 $\lambda|J_{\dR}=0$，
不是仅有某个积分 $\lambda(c(\sigma))=0$。

有限长度由忠实正合的有限维实现保证。定义
\begin{equation}\label{eq:basic}
 R_\Sigma=\bigcap_{\Sigma\subset J_B}J,
 \qquad I_\Omega=\sum_{\Omega|J_{\dR}=0}J,
 \qquad \Vis(E;\Sigma,\Omega)=R_\Sigma/(R_\Sigma\cap I_\Omega).
\end{equation}
交和求和均在实际子对象中进行。它们分别是链的范畴闭包和读出的最大不可见子对象。
当 $\Sigma,\Omega$ 各为一维时，\eqref{eq:basic} 是最小周期对象的同调描述。

\section{实际操作图表的相容最小化}
\subsection{图表、输入与所有后续观察}
设 $D$ 为有限有向图表，允许有环。每个顶点 $v$ 带对象 $E_v\in\mathcal C$，
每条箭头 $a:v\to w$ 带\emph{已经证明实际存在}的态射 $f_a:E_v\to E_w$。
设输入链空间为 $\Sigma_v\subset(E_v)_B$，
局部读出空间为 $\Omega_v\subset(E_v)_{\dR}^{\vee}$。
空路径表示恒等态射；路径 $p:v\to w$ 的复合记为 $f_p$。

一个输入在 $v$ 处现在不可见，并不保证它经 $v\to w$ 后仍不可见。
因此应同时保留全部路径读出，而不是先逐顶点最小化再任意加入操作。
令 $R_v^0=R_{\Sigma_v}$，$I_v^0=I_{\Omega_v}$，定义
\begin{align}
 R_v^{r+1}&=R_v^r+\sum_{a:u\to v}f_a(R_u^r),\label{eq:Riter}\\
 I_v^{r+1}&=I_v^0\cap\bigcap_{a:v\to w}f_a^{-1}(I_w^r).
 \label{eq:Iiter}
\end{align}
无输入时 $R_v^0=0$；无局部读出时 $I_v^0=E_v$。
逆像、交与像在 $\mathcal C$ 中均有定义。

\begin{theorem}[有限实际闭包]\label{thm:network}
令 $d=\sum_v\dim(E_v)_B$。两列 \eqref{eq:Riter}--\eqref{eq:Iiter}
分别递增、递降，至多各经历 $d$ 次严格变化后稳定。
稳定值满足
\begin{align}
 R_v&=\sum_{p:u\to v}f_p(R_u^0),\label{eq:Rpaths}\\
 I_v&=\bigcap_{p:v\to w}f_p^{-1}(I_w^0).\label{eq:Ipaths}
\end{align}
因此 $K_v=R_v\cap I_v$ 构成子图表，且
\begin{equation}\label{eq:networkQ}
 Q_v=R_v/K_v
\end{equation}
及所有诱导箭头均为实际对象、实际态射。
原输入沿任意路径的所有标架形式周期均保留。
该约化图表由输入在范畴中生成，且不存在被全部局部读出消去的非零子图表。
\end{theorem}
\begin{proof}
对路径长度归纳，$R_v^r$ 是长度不超过 $r$ 的路径像之和；
$I_v^r$ 是长度不超过 $r$ 的路径逆像之交。
于是两列分别单调。一次严格变化使实现的总维数至少增加或减少一，
而维数介于 $0$ 与 $d$ 之间。
某一步若所有顶点均不变，递推下一步仍相同，故不会在稳定之后重新增长。
这同时证明有限停止和 \eqref{eq:Rpaths}--\eqref{eq:Ipaths}。

路径复合给 $f_a(R_v)\subset R_w$ 及 $f_a(I_v)\subset I_w$，
所以 $f_a(K_v)\subset K_w$，能够形成实际商箭头。
输入属于 $R_v$，而每一后续读出均消去 $K_v$。
形式周期的函子性遂给所有路径标架的相等，不需周期求值映射单射。

若 $Q$ 的一个子图表包含所有约化输入，其在 $R$ 中的逆像仍包含原输入。
由 $R$ 的最小性，该逆像就是 $R$，故该子图表为 $Q$。
若 $L\subset Q$ 是全部局部读出不可见的子图表，其在 $R$ 中的逆像
沿任一路径仍不可见，因此属于 $I$；它已经被 $K=R\cap I$ 商去，故 $L=0$。
\end{proof}

该定理给的是有限长度构造，\emph{不是}从边表求 Nori 子对象的现成有效算法。
初始的 $R_v^0,I_v^0$ 已要求知道实际范畴闭包。
即使停止界是 $d$，若这些子对象或实际箭头未知，
也不能将未知步骤写成已完成的线性代数。

\subsection{为什么增加操作后不能沿用旧商}
设已得到 $Q_v=R_v/K_v$，现在拟加入 $f:E_v\to E_w$。
即使 $f$ 是实际态射，原商也只有在
\begin{equation}\label{eq:descent}
 f(R_v)\subset R_w,\qquad f(K_v)\subset K_w
\end{equation}
时才承载这条新箭头。第一条件防止产生旧可达对象之外的新状态；
第二条件防止把已经删掉的状态变成可观察状态。

最小的线性反例是在 $k^2$ 上只给恒等操作、输入和读出均选第一坐标。
旧最小空间为一维。加入交换两个坐标的矩阵后，
第二坐标先被访问、再返回第一坐标，旧商不再支持全部响应。
这只是操作表示反例；它并不把该交换矩阵宣称为某个未知动机的实际态射。

\begin{corollary}[物理关系的允许用法]
若分子关系、边界关系或对称恒等式已被实现为上述实际图表中的态射关系，
则它们在 \eqref{eq:networkQ} 中保持。
若只有一个数值响应为零，或者只有一条路径上的函数恒等式，
尚不足以在整个图表上商去相应方向。
\end{corollary}
\begin{proof}
实际态射关系经过子对象限制和商对象仍成立。
后一个断言由定义即可看出：$I_v$ 量化所有后续路径，
一个指定响应的消失没有提供这些条件。
\end{proof}

\section{从 Betti 操作到动机态射：商上的精确判据}
\subsection{必须检验的群是实际群}
在包含相关对象的中性 Tannakian 子范畴内，令
$G=\operatorname{Aut}^{\otimes}(\omega_B)$。
给定实际子商 $Q_i=R_i/K_i$ 及 $Q_j=R_j/K_j$，
设一个 $k$-线性 Betti 操作 $L:(E_i)_B\to(E_j)_B$ 已满足
\eqref{eq:descent}。它诱导向量空间映射 $\bar L:(Q_i)_B\to(Q_j)_B$。
这里暂不假设 $L$ 是环境对象上的动机态射。

\begin{theorem}[余模下降判据]\label{thm:intertwine}
上述 $\bar L$ 来自唯一的实际态射 $Q_i\to Q_j$，当且仅当
\begin{equation}\label{eq:commutator}
 \bigl(\rho_j(g)L-L\rho_i(g)\bigr)(R_i\otimes A)
       \subset K_j\otimes A
\end{equation}
对每个 $k$-代数 $A$ 及每个 $g\in G(A)$ 成立。
等价地，设 $q_j:(R_j)_B\to(Q_j)_B$，以
$\delta_i:(R_i)_B\to(R_i)_B\otimes\mathcal O(G)$ 表示余作用，则
\begin{equation}\label{eq:comodule}
 (q_j\otimes\id)\delta_j L|_{R_i}
       =(q_jL|_{R_i}\otimes\id)\delta_i.
\end{equation}
给定一个预期的 de Rham 映射时，还必须核对它与这个唯一实际态射的
de Rham 实现一致。
\end{theorem}
\begin{proof}
由于 $R_i,K_i,R_j,K_j$ 为实际子对象，它们的 Betti 实现对 $G$ 稳定。
因此 \eqref{eq:commutator} 恰好说 $\bar L$ 是 $G$-等变映射，
\eqref{eq:comodule} 是同一条件的坐标环表述。
中性 Tannaka 等价的满性给实际态射，忠实性给唯一性。
反向则由实际态射对纤维函子所有张量自同构的自然性成立。
\end{proof}

这个判据有两点作用。其一，$L$ 可以在环境上不等变，而在商去不可见核后等变；
所以要求环境交换子严格为零可能过强。其二，不能把 $G$ 偷换为容易计算的
单值子群或某一组形式操作。对较小群等变一般只给必要条件。
若实际群尚未知，定理保留为待检验的条件，不声称提供任意图的自动提升。

在有限维表示上，\eqref{eq:comodule} 的两边只使用有限个坐标函数，
因此\emph{在实际余作用已经给出时}可变成有限线性条件。
这不要求仅用 $G(k)$，也不默认有理点 Zariski 稠密。

\subsection{自同态的直接障碍}
\begin{corollary}\label{cor:scalar}
若 $R=E,K=0$ 且 $\End_{\mathcal C}(E)=k$，
则能成为该固定最小对象自同态的 Betti 操作恰为标量矩阵。
特别地，一个非标量单值反射不能由``它保持配对''推出为该对象的动机自同态。
\end{corollary}
\begin{proof}
定理\ref{thm:intertwine}将允许映射识别为
$\End_G(E_B)=\End_{\mathcal C}(E)$，即标量。
保持一个配对只表明矩阵属于一个等距群，不表明它与实际 Galois 作用交换。
\end{proof}

真正的割操作并不因此被排除。留数的对偶管态射通常是
$f:C\to E$，其\emph{源为另一个带 Tate 扭转的实际对象}；
它不是任意给定的 $E\to E$ 矩阵。
应检验 $f$ 的实际来源、链标记以及 $f^*\lambda$，而不是把局部单值本身当成 $f$。

\section{head 与 socle：保护和最小性的不同问题}
\subsection{单源保护只约束非零商}
在有限长度范畴中，$\head R$ 表示最大半单商，
$\soc R$ 表示所有简单子对象之和。对固定有限对象可转到其有限维表示代数，
二者均存在。

\begin{proposition}[范畴版单源 head 判据]\label{prop:head}
设 $R=R_\sigma$，$S$ 为简单对象。
每个保留 $\sigma$ 的非零商都含有 $S$ 作为组成因子，
当且仅当 $\head R$ 的全部简单因子都同构于 $S$。
\end{proposition}
\begin{proof}
一个非零商中的 $\sigma$ 不可能为零，否则商映射的核承载 $\sigma$，
与 $R=R_\sigma$ 矛盾。非零有限长度商有简单商；
它也是 $R$ 的简单商，所以正向保护由 head 的单类型条件得到。
若 head 含 $S'\not\simeq S$，投到 $S'$ 即给出不含 $S$ 的非零商。
\end{proof}

这里量化所有实际商，尚未证明每个商都可由允许物理分子实现。
因此这是充分保护判据及所述\emph{逻辑量化}下的必要条件，
不是任意指定物理读出族的无条件充要分类。

\subsection{读出是否完整，由简单子对象检测}
\begin{proposition}[socle 检测]\label{prop:socle}
设 $R=R_\Sigma$。最小化不再损失子对象，即
$\Vis(R;\Sigma,\Omega)=R$，当且仅当不存在非零简单子对象
$S\subset R$ 被全部读出消去。
特别地，若 $\soc R=A$ 是简单对象，则
\begin{equation}\label{eq:socle}
 \Vis(R;\Sigma,\Omega)=R
 \quad\Longleftrightarrow\quad \Omega|A_{\dR}\ne0.
\end{equation}
\end{proposition}
\begin{proof}
任一非零不可见子对象在有限长度范畴中包含简单子对象。
反之，一个被消去的简单子对象已经属于最大不可见核。
若 socle 简单，则每个非零子对象都含 $A$，所以 \eqref{eq:socle} 立即成立。
\end{proof}

非半单对象中，简单纯核可以处于 socle，而顶层 Tate 商处于 head。
``一个输出非零''与``输出在纯核上的限制非零''此时强度完全不同。
这个区别正是下面 banana 应用的关键。

\section{指定 banana 对象的全部子对象}
\subsection{已建立的几何输入与适用范围}
本节采用已有证明稿给出的实际同调对象 $E$，基域为 $\Q$，
固定 $m=(1,2,4,8)$、$s=-64$。它具有实际过滤
\begin{equation}\label{eq:banana}
 0\subset A\subset N\subset E,
 \qquad \dim A=6,\quad N/A=B=\Q(1),\quad E/N=\mathbf1.
\end{equation}
$A$ 是相应扭转的简单 K3 超越对象。
已有相邻扩张非零证明和实际群论证给出完整幺幂根基
\begin{equation}\label{eq:U}
 U=\left\{
 \begin{pmatrix}I_A&x&z\\0&1&y\\0&0&1\end{pmatrix}
 :x\in\Hom(B,A),\ y\in B,\ z\in A\right\}.
\end{equation}
这是在一个 Levi 分裂中显示的式子；$A,N,E$ 本身不依赖该分裂。
纯分级作用在 $A$ 上不可约。下述证明只需要这些输入，
不再使用新的 Frobenius、周期或路径计算，也不需要重新判别固定纯群的分支。

指定点的输入分别见内部证明稿
\nolinkurl{LOWER_REGULATOR}、\nolinkurl{MIXED_UNIPOTENT_GROUP}、
\nolinkurl{FIXED_PURE_GROUP} 及主稿第6--7、10节。
一般三步对象若缺少 \eqref{eq:U} 或纯核简单性，不能套用以下分类。

\begin{theorem}[全部实际子对象]\label{thm:uniserial}
对象 $E$ 的实际子对象恰好为
\begin{equation}\label{eq:chain}
                        0,\ A,\ N,\ E.
\end{equation}
因此 $\soc E=A$，$\head E=\mathbf1$。
\end{theorem}
\begin{proof}
令 $J\subset E$ 非零，在 \eqref{eq:U} 的分裂中检查其 Betti 实现。
若 $J$ 在顶层非零，取 $v=(a,b,c)\in J_B$ 且 $c\ne0$。
仅改变 $z$ 的群元素给出 $(u_z-1)v=(cz,0,0)$，故 $A_B\subset J_B$。
仅改变 $y$ 给 $(u_y-1)v=(0,cy,0)$，故整个中间 $B$ 方向也在 $J_B$ 中。
于是 $N_B\subset J_B$，加上原顶层向量即得 $J=E$。

若顶层为零但 $J$ 在 $N/A$ 中非零，取 $v=(a,b,0)$ 且 $b\ne0$。
由全部 $x$ 得 $(u_x-1)v=(xb,0,0)$，所以 $A_B\subset J_B$，
再加该中间向量得 $J=N$。
最后，若 $J\subset A$，则纯核的简单性给 $J=A$。
上述包含通过忠实正合实现反映回实际范畴，四个候选均已实际存在。
\end{proof}

这个证明没有把不可分解等同于简单。
$E$ 有非平凡子对象 $A,N$，但没有额外直和投影。
由 \eqref{eq:chain} 和扩张的非分裂，非零动机自同态必须同时尊重各层；
已有群作用进一步给 $\End(E)=\Q$。

\section{完整双标架分类与分子选择规则}
\subsection{链的闭包与读出的核}
对任意 $\sigma\in E_B$，定理\ref{thm:uniserial}立即给
\begin{equation}\label{eq:Rlevels}
 R_\sigma=
 \begin{cases}
 0,&\sigma=0,\\
 A,&0\ne\sigma\in A_B,\\
 N,&\sigma\in N_B\setminus A_B,\\
 E,&\sigma\notin N_B.
 \end{cases}
\end{equation}
对任意 $\lambda\in E_{\dR}^{\vee}$，有
\begin{equation}\label{eq:Ilevels}
 I_\lambda=
 \begin{cases}
 0,&\lambda|A_{\dR}\ne0,\\
 A,&\lambda|A_{\dR}=0,\ \lambda|N_{\dR}\ne0,\\
 N,&\lambda|N_{\dR}=0,\ \lambda\ne0,\\
 E,&\lambda=0.
 \end{cases}
\end{equation}
这里均指向量或限制映射的精确非零，不是十进制求值不为零。

\begin{theorem}[所有双标架的最小对象]\label{thm:allframes}
记 $K=E/A$，它是非分裂的秩二质量 Kummer 对象。
非零输入的完整分类为下表；括号中为秩。
\begin{center}
\small
\begin{tabular}{c|cccc}
\toprule
链的范畴闭包&$\lambda|A\ne0$
 &$\lambda|A=0,\lambda|N\ne0$
 &$\lambda|N=0,\lambda\ne0$&$\lambda=0$\\
\midrule
$A$&$A\ (6)$&$0$&$0$&$0$\\
$N$&$N\ (7)$&$B\ (1)$&$0$&$0$\\
$E$&$E\ (8)$&$K\ (2)$&$\mathbf1\ (1)$&$0$\\
\bottomrule
\end{tabular}
\end{center}
各项都是实际最小对象；在通常形式周期记法中取它的对偶。
\end{theorem}
\begin{proof}
将 \eqref{eq:Rlevels}、\eqref{eq:Ilevels} 代入
$R_\sigma/(R_\sigma\cap I_\lambda)$。
对象识别分别使用 $N/A=B$、$E/A=K$ 和 $E/N=\mathbf1$。
没有剩余未知子对象，故该表既给上界也给最小性。
\end{proof}

\subsection{固定物理链的精确秩分层}
原正物理链 $\sigma_{\mathrm{phys}}$ 在 $E/N$ 的坐标已标定为 $1$，
所以其范畴闭包为 $E$。
现在固定它，仅允许改变实际 de Rham 读出。
\begin{corollary}[固定链的四个层]\label{cor:numerators}
设 $V$ 为一个有限维 $\Q$-空间，代表已经映入同一对象的允许分子类；
设线性映射 $L:V\to E_{\dR}^{\vee}$ 给其读出。
定义
\begin{equation}\label{eq:exceptional}
 Z_A=\ker(\operatorname{res}_A L),\qquad
 Z_N=\ker(\operatorname{res}_N L),\qquad Z_E=\ker L.
\end{equation}
则 $Z_E\subset Z_N\subset Z_A\subset V$，且
\begin{equation}\label{eq:rankstrata}
 \rank\Vis(E;\sigma_{\mathrm{phys}},L(v))=
 \begin{cases}
 8,&v\notin Z_A,\\
 2,&v\in Z_A\setminus Z_N,\\
 1,&v\in Z_N\setminus Z_E,\\
 0,&v\in Z_E.
 \end{cases}
\end{equation}
某些层可以为空。该分类不要求先求出任一个完整积分。
\end{corollary}
\begin{proof}
直接应用定理\ref{thm:allframes}的最后一行。
每个例外集合都是限制映射的核，故是有理线性子空间。
若允许读出为空间 $E_{\dR}^{\vee}$ 本身，则
$\dim Z_A=2,\dim Z_N=1,\dim Z_E=0$；
对于物理给定子空间，只取这些空间的逆像，不能默认其全满。
\end{proof}

本推论的前提是分子已给出\emph{同一相对对象}上的类。
若分子增加极点除子、要求新增 sector，或改变边界及链，须先扩大或重构对象。
尤其不能仅按多项式次数将任意分子自动放入这里的八维空间。
多分子线性组合的相消已经由 $L$ 的核精确记录；
积分后一个值恰为零，并不意味着 $v\in Z_E$。

\subsection{K3 可见为何迫使混合扩张一起保留}
在保留原链时，$\lambda|A\ne0$ 不仅保护 K3 简单因子，
还使整个 $E$ 最小。这是因为任一非零子对象都包含 $A$。
不能先商去中间层而保留一个同样带原链的纯 K3 商；
\eqref{eq:chain}中没有这种商。

相反，最大割把链变成 $A_B$ 中的非零向量，并使用相应非零纯读出。
定理\ref{thm:allframes}第一行给纯秩六对象。
下层周期取 $N_B\setminus A_B$ 的链，则可能得到秩七。
因此 $6$、$7$、$8$ 的区别来自完整双标架，不能归结为环境中是否存在一个 K3 曲面。

同时，$\head E=\mathbf1$ 表明：仅凭一个未知读出的非零性，
在实际动机范畴中只可强制顶层简单类型，不能单独强制 K3。
顶层商读出 $E\to\mathbf1$ 已给非零但不含 K3 的反例。
原切割单值模的 head 为 K3 型时仍可使用保护定理；
那是另一个对象、另一组链和一个不同的作用范畴，不存在矛盾。

\section{幺幂深度与实际割见证}
\subsection{三个非零层的深度}
\begin{proposition}\label{prop:depth}
对固定顶层非零链，\eqref{eq:rankstrata}的三个非零层分别具有
形式周期的幺幂深度 $2,1,0$。
\end{proposition}
\begin{proof}
在 $E$ 上，$u-1$ 严格上三角，所以任意三个这样的因子的乘积为零。
当 $\lambda|A\ne0$ 时，选 $x,y$ 使 $\lambda(xy\sigma_{\mathrm{top}})\ne0$；
这是可能的，因为 $\sigma_{\mathrm{top}}\ne0$ 且
$\Hom(B,A)\otimes B\to A$ 的评价满射。
于是一个二重幺幂变化非零，深度恰为二。
此计算在比较扭子的任意平凡化上进行，故是形式余作用的断言。

在商 $K=E/A$ 上，幺幂作用只剩 $y$，其二重变化全为零。
读出在 $B$ 上非零且链顶层非零，故一重变化非零，深度一。
再商为 $\mathbf1$ 后幺幂作用平凡，深度零。
\end{proof}

深度不是数值超越性标签。秩六的纯 K3 对象同样深度零，
而秩一的 Tate 周期也深度零。深度计量扩张中的幺幂方向，
纯核类型必须单独保存。本文也不由深度二推断函数是某种长度二的 MPL。

\subsection{用真实割证明读出在纯核上非零}
\begin{proposition}[实际割的充分见证]\label{prop:cutwitness}
设一个实际态射 $f:C\to E$ 的像恰为 $A$，
并有实际割链 $\gamma\in C_B$。
若 $\langle f^*\lambda,\gamma\rangle\ne0$，
则 $\lambda|A\ne0$。
若原链在 $E/N$ 非零，则原标架的最小对象为完整 $E$。
\end{proposition}
\begin{proof}
若 $\lambda|A=0$，则 $f^*\lambda=0$，所有该类割周期都为零，矛盾。
再用定理\ref{thm:allframes}即可。
从一个实际数值\emph{非零}推出形式或向量非零不需要周期猜想；
危险的是反方向，把数值零当成形式零。
\end{proof}

必须证明 $f$ 的像为纯核，而不能仅因某个正的相空间积分非零就假定这一点。
若割对象还含可见边界项，非零也可能完全来自该边界。
此外，函数在一个参数域内不恒零，不保证在任一指定代数点仍非零。
对指定点使用命题\ref{prop:cutwitness}时，应提交该点的见证，
或直接以代数 de Rham 限制证明 $\lambda|A\ne0$。

对于本项目已给定的原物理读出，该限制非零来自既有全纯形式与纯核比较，
不需要为本命题再做一次数值周期计算。新分子仍须提交其实际类及限制。

\section{可复用规则及未完成范围}
本补篇给统一分类方案增加三个可直接调用的理论规则。

第一，操作闭包必须以实际图表为输入，并同时考虑所有后继读出。
定理\ref{thm:network}给规范约化图表；
将一个算子加入既有商时，先检验可达子对象和不可见核是否保持。

第二，已下降的线性操作仍须通过实际余作用的等变检验。
定理\ref{thm:intertwine}将这个责任精确地写在商上，
既不误要求环境等变，也不把单值、配对保持或 Hodge 相容直接视为动机性。

第三，在已识别简单 socle 的实际对象中，
纯核上的非零读出可比单纯的 head 保护强得多。
指定 banana 的完整表给出全部双标架最小对象及固定链分子分层，
其中混合扩张的共同保留来自实际子对象结构。

这些规则没有构造尚缺失的一般图 Nori 群、所有规范投影、
任意分子的统一相对模型或每个物理商的可实现性。
正质量零动量森林的实际对象可以作为图表顶点，
但减除或局部 RG 提取若要成为箭头，仍需各自的几何比较。
当前工作也不将零点森林的已有收敛证明视为本批重新审计完成。

后续理论工作可先选一个具有明确相对分子类的有限族，
在同一对象中写出 $L$ 及限制映射，得到具体线性消零条件；
或者为一个实际减除/局部投影补出图表态射。
这两条路线都以几何定义及证明为主体，不以扩大数值矩阵作为替代。

\begin{thebibliography}{9}
\bibitem{brown}F. Brown,
\emph{Notes on motivic periods}, arXiv:1512.06410v3 (2017),
\S\S2.4--2.5.
\url{https://arxiv.org/html/1512.06410}.
最小对象和 coradical 为原文基础；本文的图表及标架应用另证。

\bibitem{dm}P. Deligne and J. S. Milne,
\emph{Tannakian Categories}, in \emph{Hodge Cycles, Motives, and Shimura Varieties},
Lecture Notes in Mathematics 900, Springer (1982), pp. 101--228.
使用作者维护版本 Theorem 2.11：
\url{https://www.jmilne.org/math/xnotes/tc.pdf}.

\bibitem{internal}费曼周期研究项目，
\emph{protected\_period\_network}、\emph{framed\_exact\_classification}（2026年10月10日），
以及 \emph{VISIBILITY\_AND\_CORADICAL}、\emph{THEORY\_MANUSCRIPT} 第6--7、10节。
内部工作稿；定理输入范围及有限证据由相应证明稿记录，不表示外部同行评议已完成。
\end{thebibliography}
\end{document}
